\emph{By James Radek}

According to the dominant view on the broad liberal-left, to achieve progressive change we must elect officials with progressive politics. Liberalism claims we can eliminate the problems inherent in capitalism without eliminating capitalism itself, and by using the system created by capitalists to maintain their rule.

Social democracy (or ``democratic socialism'') takes this na\"{\i}ve error even further. Social democrats believe we can eliminate capitalism itself within the legal and political framework of capitalism. No systemic rupture is required---we simply need to elect a critical mass of socialists to legislate a critical mass of reforms.

Revolutionary socialism holds that we must fight for reforms, but the most important part is \emph{the fight}---and the organization required to mount it---not the reforms themselves.

We view all popular grievances---like poverty, war, ecological collapse, racial, national, sexual, and gender oppression---as interconnected, sharing a common source: capitalism. So the goal of Marxists is to change the global mode of production, the way we make things, from capitalism to socialism.

Marxists are champions of democracy. But only a highly circumscribed, sham democracy is possible under capitalism, where political power flows solely from economic power.

Marxists understand the state to be, in Engels's words, ``nothing but a machine for the oppression of one class by another, and indeed in the democratic republic no less than in the monarchy.''\footnote{Frederick Engels, ``Introduction'' to \emph{The Civil War in France}.} And ``the executive of the modern state,'' according to \emph{The Communist Manifesto}, ``is but a committee for managing the common affairs of the whole bourgeoisie.''

\blockquote{Our view of the capitalist state informs the strategies and tactics of Marxists. We see it as a weapon in the hands of the capitalist class, an antagonistic force that must be fought and ultimately smashed, not gradually improved, strengthened, or integrated with by winning office.}

This may be more apparent when it comes to presidents and congresspeople, but local and regional politics offer no alternative. Every mayor is a servant of real-estate developers and whatever corporate enterprises dominate their city. The job of governors in states like Colorado is to manage the extraction of fossil fuels.

We are in the minority with this view, even among socialists. The Democratic Socialists of America (DSA) is directing radicals into an imperialist party. We're trying to do the opposite. As Trotsky said, ``Social Democracy considers itself the democratic doctor of capitalism, we are its revolutionary gravediggers.''\footnote{Leon Trotsky, ``The Tragedy of the German Proletariat,'' 1933.} Social-democratic ideas are providing grave misleadership to the current socialist movement by redirecting discontent, grassroots struggles, and the rising interest in socialism into the Democratic Party and electoral politics more broadly.

\paragraph{The Democratic Party is a capitalist and imperialist party}

The Democratic Party is an irredeemably undemocratic tool of the capitalist class that cannot be reformed, utilized, or taken over by the working class and oppressed. On the contrary, it remains an impediment to working-class consciousness and organization, \emph{not in spite of, but due to} the growth of a so-called socialist wing within it. The Democrats' historic role as ``the graveyard of social movements'' has not changed. Its role is to absorb and neutralize discontent by rerouting it into an arena controlled by capital, where it can do no harm.

As the parties of the liberal bourgeoisie have done throughout history, the Democrats tend to the longer-term interests of the system by making minor concessions and major rhetorical appeals to our interests. Franklin Delano Roosevelt, for instance, was class-conscious of his role in granting the New Deal, which staved off revolts and revolution, while advancing the US imperial project and maintaining the racist regime of capital. As he once said, ``I am the best friend the profit system ever made.''\footnote{Quoted in \emph{Socialist Worker}, November 14, 2008.}

Whether conscious of it or not, contemporary left Democrats play a similar role in disorienting would-be militants. Bernie Sanders, the Squad, and other democratic socialists didn't popularize socialism. They gained popularity by associating themselves with socialism. And they distorted its meaning beyond recognition in the process. Now loads of random liberals are declaring themselves socialists just to get the DSA to campaign for them.

\paragraph{The urgency of class independence}

While most of the liberal-left in the US seeks advance through the Democratic Party and the electoral system more broadly, Marxists point to class struggle as the motor force of history. Specifically this means mass, independent, militant, working-class action. Independent means \emph{class-}independent---independent of the big and petty bourgeoisies. It means not sharing organizations with them, and not taking money from them, as does the Democratic Party and the NGO-industrial complex that surrounds it.

Class struggle is primarily advanced through strikes, protests, occupations, and uprisings, and the organizational forms required to cohere the class and take such actions, like trade unions, tenants' unions, workers' councils, and revolutionary organizations. It must not be limited by what is legal, as the law flows from and protects the profit system. Everything we do in the here and now should be in the service of the future emergence of such class coherence, organization, and combativity.

We cannot advance toward revolution unless an increasing number of workers engage in mass, independent, militant action and organize themselves accordingly. Integral to this process are: political education and experience in struggle---struggle against our enemies \emph{and} amongst our allies.

\paragraph{How is progressive change achieved under bourgeois democracy?}

Social progress occurs when the working class and oppressed make temporary advances in their irreconcilable conflict with the capitalist class. All major social reforms in US history have been the product of mass, militant struggle: from labor rights, unemployment compensation, welfare, and overtime pay in the 1930s to civil rights, abortion rights, environmental protection, affirmative action, and Medicare in the 1960s and 70s.

As Trotsky wrote, ``The general Marxist thesis, `Social reforms are only the by-products of the revolutionary struggle,' has in the epoch of the decline of capitalism the most immediate and burning importance. The capitalists are able to cede something to the workers only if they are threatened with the danger of losing everything.''\footnote{Leon Trotsky, \emph{Whither France?}, 1936.} By contrast, in Lenin's words, ``Reformist tactics are the least likely to secure real reforms.''\footnote{V. I. Lenin, ``The Question of a National Survey of the Agrarian Question,'' 1906.}

\blockquote{As the parties of the liberal bourgeoisie have done throughout history, the Democrats tend to the longer-term interests of the system by making minor concessions and major rhetorical appeals to our interests. Franklin Delano Roosevelt, for instance, was class-conscious of his role in granting the New Deal, which staved off revolts and revolution, while advancing the US imperial project and maintaining the racist regime of capital.}

A ruling class may grant reforms when threatened, but it may also be willing to dispense with the fa\c{c}ade of democracy altogether when necessary. ``The constitutional road is cut by trenches held by fascist bands,'' as Trotsky wrote.\footnote{Leon Trotsky, \emph{Whither France?}, 1936.}

The 1973 coup against Salvador Allende's democratically-elected socialist government in Chile is a case in point.\footnote{See Peter Gonzalez, ``Chile 1973: Was Victory Possible?''} As any ruling class around the world would do when threatened by a rise in class struggle, Chilean capitalists dropped the democratic fa\c{c}ade, slaughtered and jailed thousands, erased the gains made under Allende and before, and restored the economy to its role as their profit-making machine. The working class movement was systematically destroyed by the Pinochet regime, perhaps most dramatically via the mass internment, torture, and execution of its suspected political enemies at the National Stadium of Santiago.

The electoral road to socialism ends in a Chilean soccer stadium. Reformist socialists have not learned this lesson.

\paragraph{What can we learn from the history of the socialist left's approach to elections?}

The Marxist political tradition warns us against what Engels called ``parliamentary cretinism''---the tendency to place all hope for social change in parliamentary institutions.\footnote{Frederick Engels, ``Germany: Revolution and Counter-Revolution,'' 1852.} But even socialists trained in Marx have shown they are not immune.

Before 1900, when many dozens of German socialists won election to parliament, Engels became increasingly alarmed at the rise of reformist tendencies. Karl Kautsky, leading member of the German Social Democratic Party (SDP) and once known as the ``Pope of Marxism,'' went so far as to call parliamentary activity ``the most powerful lever that can be utilized to raise the proletariat out of its economic, social and moral degradation.''

At the turn of the century, Rosa Luxemburg railed against such tendencies:

\blockquote{People who pronounce themselves in favor of the method of legislative reform \emph{in place and in contradistinction to} the conquest of political power and social revolution, do not really choose a more tranquil, calmer and slower road to the \emph{same} goal, but a \emph{different} goal. Instead of taking a stand for the establishment of a new society they take a stand for surface modifications of the old society.\footnote{Rosa Luxemburg, \emph{Reform or Revolution}, 1900.}}

The history of socialists' attempts to reform the state and vote socialism into power has resulted in socialists adapting to nationalism, becoming managers of austerity and imperialism, disorienting and disorganizing the working class, and making no advance toward overturning the reign of capital. This is true across the history of European socialism---from the Second International's betrayal of the international working class by supporting ``their own'' national bourgeoisies in World War I, to the eurocommunism of the 1970s and 80s, all the way up to SYRIZA in 2015.\footnote{See V. I. Lenin, ``The Tasks of the Proletariat in Our Revolution,'' 1917.}

\blockquote{We cannot advance toward revolution unless an increasing number of workers engage in mass, independent, militant action and organize themselves accordingly. Integral to this process are: political education and experience in struggle---struggle against our enemies \emph{and} amongst our allies.}

In the 1930s, under the influence of Stalin, the Communist Party in the US pursued a class-collaborationist policy that led them to support FDR and the Democrats.\footnote{See ``The US Communist Party's Class Collaborationist Past,'' \emph{Red Flag}.} This undermined the possibilities for a militant workers' movement, significant reforms, the popularizing of left-wing politics, or the emergence of an independent labor party in the United States. Today, the DSA is doing the same thing, squandering the rising interest in socialism.

Many other US socialists in the early to mid-20th century, including James Cannon, Hal Draper, and Eugene Debs, upheld the Marxist principle of class independence.

In 1948, James Cannon offered this defense of class independence that is worth quoting at length:

\blockquote{It has been argued here that `we must go through the experiences with the workers.' That is a very good formula, provided you do not make it universal. We go with the workers only through those experiences which have a class nature. We go with them through the experiences of strikes, even though we may think a given strike untimely. We may even go with the workers through the experience of putting a reformist labor party in office, provided it is a real labor party and subject to certain pressures of the workers, in order that they may learn from their experience that reformism is not the correct program for the working class.

But we do not go through the experience of class collaboration with the workers. There we draw the line. We did not go through the experience of the workers when they supported the imperialist war. We drew back when they went through the experience of people's fronts in Europe. We stood on the side and we told them they were wrong. We did not compromise ourselves. If another man takes poison, you do not have to join him in the experiment. Just tell him it is no good. But don't offer to prove it by your personal example.\footnote{James Cannon, ``Defense of Class Independence,'' 1948.}}

\paragraph{What does a revolutionary socialist approach to elections look like?}

While Marxists do not consider the capitalist state to be a vehicle for socialist transformation, we also oppose abstentionism. We don't counterpose elections to strikes, protests, occupations, and uprisings. Rather, we seek to integrate all forms of struggle into a comprehensive strategy for revolution. While we decry the limitations of bourgeois democracy, we prefer it to the alternatives under capitalism, such as authoritarianism and fascism. Political freedoms, like the right of assembly, association, free speech, and a free press, are valuable gains of past struggle that help us further our own class interests.

For Marx and Engels, electoral victories were subordinate to independent working-class political action. As an example, in 1850, they gave an address to the Communist League in London that, 50 years later, Lenin found so compelling that he memorized it, enjoyed reciting it, and fought for it to become the core of Bolshevik electoral strategy.\footnote{Karl Marx and Frederick Engels, ``Address of the Central Committee to the Communist League,'' 1850.} Here's a key passage from that address:

\blockquote{Even when there is no prospect whatever of their being elected, the workers must put up their own candidates in order to preserve their independence, to count their forces and to lay before the public their revolutionary attitude and party standpoint. In this connection they must not allow themselves to be bribed by such arguments of the democrats as, for example, that by doing so they are splitting the democratic party and giving the reactionaries the possibility of victory. The ultimate purpose of all such phrases is to dupe the proletariat. The advance which the proletariat party is bound to make by such independent action is infinitely more important than the disadvantage that might be incurred by the presence of a few reactionaries in the representative body.}

To break this down, first of all: winning isn't everything! The point of a revolutionary socialist electoral campaign is \emph{not} to get elected. We don't want to become managers of austerity and imperialism for the capitalist state. We want to overthrow it. We don't want to dedicate our lives to tweaking the system or winning tiny concessions that can be swept away with the stroke of a pen by the next elected official.

\blockquote{The history of socialists' attempts to reform the state and vote socialism into power has resulted in socialists adapting to nationalism, becoming managers of austerity and imperialism, disorienting and disorganizing the working class, and making no advance toward overturning the reign of capital.}

We do want to ``lay before the public our revolutionary attitude and party standpoint''---in other words to propagate socialist politics and build revolutionary organization. Electoral campaigns and elected offices can serve as big platforms to expose the system and encourage struggle against it.

And the insistence on class independence, even when it may lead to right-wing victories, is another important position that is sadly unpopular on the liberal-left, where safe-state strategies and lesser-evilism dominate. In 2020, for instance, the Party for Socialism \& Liberation (PSL) ran its own candidate for president---but if Bernie Sanders had won the Democratic nomination, they committed to not running their candidate in swing states to avoid splitting the vote.

In 1906, Lenin addressed the matter of socialists forming blocs with liberals to stave off reactionaries in this way:

\blockquote{We would undermine the principles and the general revolutionary significance of our campaign for the sake of gaining a seat of the Duma for a liberal! We would be subordinating class policy to parliamentarism instead of subordinating parliamentarism to class policy. We would deprive ourselves of the opportunity to gain an estimate of our forces. We would lose what is lasting and durable in all elections---the development of the class consciousness and solidarity of the socialist proletariat. We would gain what is transient, relative and untrue---superiority of the Cadet over the Octobrist.\footnote{V. I. Lenin, ``The Elections to the Duma and the Tactics of the Russian Social-Democrats,'' 1906.}}

Nowhere in Marx and Engels's 1850 address do they mention changing the system from within. The focus is always on the self-activity of militant workers. Lenin warned that socialists elected to the First Duma should ``avoid getting sucked into the minutia of law making''---what Engels called ``project-mongering''---and focus on making appeals to the people. This approach was later generalized by the Third International in 1920, which resolved that, ``The Communist Party does not enter these institutions [bourgeois governments] in order to carry out organic work there [reforms], but in order to help the masses from inside parliament to break up the state machine and parliament itself through action.''\footnote{Second Congress of the Communist International, ``On the Parliamentary System,'' 1920.}

For Marxists, class struggle is both the way to win reforms now and the basis for the future revolutionary transformation of capitalism into socialism. Marxists assess electoral campaigns according to how they impact class struggle. When workers break away from support for capitalist parties and run their own working-class candidates against them, this advances class consciousness and can lay the basis for an increase in class struggle. This is why Marxists agitate for the formation of a working-class political party. We do not think that electing a labor party to office within the capitalist electoral system can bring about socialism---that will require revolution.

\paragraph{Conclusion}

So, where does all of this leave us? Should we support the lesser evil? Hal Draper, then of the Independent Socialist Club, dealt with this question in the late 60s:

\blockquote{So who was really the Lesser Evil in 1964 [Lyndon Johnson or Barry Goldwater]? The point is that it is the question which is a disaster, not the answer. In setups where the choice is between one capitalist politician and another, the defeat comes in accepting the limitation to this choice.\footnote{Hal Draper, ``The Lesser of Two Evils,'' 1967.}}

We stand in that political tradition. Lesser-evilism has no bottom, no lower limit to the politics one must support by its logic, as ``Genocide Joe'' Biden has proven in recent months. It only lowers our horizons and keeps us fixated on electoral solutions.

Should we support left candidates outside of the two dominant capitalist parties? That was the practice of the International Socialist Organization, which many of us came out of. The ISO was so fixated on breaking with the Democrats that it ended up supporting and running candidates through the Green Party, a horribly confused, politically eclectic, petit-bourgeois formation. That is not what is meant by independent.

What is the point of supporting a candidate from an organization you are not a member of and do not wish to build? Just to relate to people excited by their campaign? We don't need to join a campaign to do that.

This orientation also resulted in a critique of electoralism that didn't go far enough. It seems the ISO was uncomfortable with lacking an answer on who to vote for.

\blockquote{There is a crisis of legitimacy around elections and, while most socialists try to assuage people's concerns and urge them back to the ballot box, we have chosen to affirm the crisis.}

In \emph{Prisoners of the American Dream}, the late, great Mike Davis wrote:

\blockquote{In no other capitalist country is mass political abstentionism as fully developed as in the United States, where a `silent majority' of the working class has sat out more than half the elections of the last century. Arguably, this mute, atomized protest is the historical correlative of the striking absence of an independent political party of the proletariat in the country that once invented both the labor party and May Day.\footnote{Mike Davis, \emph{Prisoners of the American Dream} (New Left Review, 1980).}}

Identifying with mass abstentionist instincts doesn't mean we have to become abstentionists though. Whether we run our own candidates is not a question of principle, but tactics. For now, we are far too small to devote the enormous energy it would take to mount electoral campaigns that would serve our goals. The electoral arena is also dangerous territory to enter for a revolutionary organization that is not unwaveringly committed to the Marxist approach I've tried to outline here. The siren song of reformism can be very tempting. And to run for office while declaring that bourgeois democracy only veils the rule of capital is a tricky needle to thread from a messaging standpoint.

When questions about what to do at the ballot box arise, revolutionary socialists must get comfortable answering \emph{there is no one to vote for}. There are no electoral solutions to the problems we face. Our promise and power lie elsewhere.
